\section{比較不等式と主定理の証明}\label{sec:conclusion}

\begin{proposition}[対数二乗余体積と指数の下界]\label{prop:lower-bound}
\(y=(y_1,\ldots,y_n)\) は \(V\) の一次独立な族で
\(n\ge(1-4/b^2)h\) とする。
その対数二乗余体積が \(c\) 以下、対数指数が \(\delta\) 以下ならば
\begin{equation}\label{eq:lower-bound}
  h(H_0-3E_4)\le c+2\delta .
\end{equation}
\end{proposition}

\begin{proof}
命題 \ref{prop:determinant-inputs} の補助形式から
\(n'=h-n\) 個を選び、\(y\) を基底に延長する。
同じ命題の \(\xi_\iota,\gamma(\xi)\) を用いて、命題 \ref{prop:gram-product}
に \(B_{\mathrm{Ar}}=2x_1^2b^2\)、\(B_{\mathrm{fin}}=b\)、
\(B_Z=h(b+E_4-H_0)\) を代入する。すると
\[
 0\le c+2\delta+n'(2x_1^2b^2+b)+h(b+E_4-H_0).
\]
補う個数は \(n'\le4h/b^2\) であり、
許容データの \(4(2x_1^2+b^{-1})\le b\) を使うと
\(n'(2x_1^2b^2+b)\le hb\) となる。
\(b\le E_4\) を用いて \(hb\) の二つの寄与を吸収すれば
\eqref{eq:lower-bound} を得る。
\leanlinks{\leanref{integral_forms_long}}
\end{proof}

\begin{theorem}[比較不等式]\label{thm:main-inequality}
任意の許容データについて
\begin{equation}\label{eq:main-ineq}
 m\mathcal R\le fa(\log b)^5 .
\end{equation}
\end{theorem}

\begin{proof}
定理 \ref{thm:spanning} の平行移動の集合 \(T\) を取る。
定理 \ref{thm:upper-bound} (1) の短い整形式が \(S\) を張るので、
それらの像 \(\mathfrak p_g\phi\), \(g\in T\) から、
\(\sum_{g\in T}\mathfrak p_g S\) の基底 \(y_1,\ldots,y_n\) を選べる。
すると \((1-4/b^2)h\le n\le h\)、かつ \(h>0\) である。
各像に対する二乗ノルムの上界と Hadamard の不等式より、
対数二乗余体積の上界として
\[
   c=n(H_0-m\mathcal R+E_4)
\]
が使える。また定理 \ref{thm:upper-bound} (2) より対数指数の上界として
\[
  \delta=2|T|(\dim S)E_3\le 10h\log(1+2b^2)E_3
\]
が使える。ここで左辺を使った指数評価を右辺へ大きくしてもよい。
命題 \ref{prop:lower-bound} から
\[
 h(H_0-3E_4)
  \le n(H_0-m\mathcal R+E_4)+20h\log(1+2b^2)E_3.
\]
\(H_0+E_4\ge0\) なので \(n(H_0+E_4)\le h(H_0+E_4)\) である。
右辺にこの評価を入れると \(hH_0\) が消え、
\[
 n\,m\mathcal R
   \le h\bigl(4E_4+20\log(1+2b^2)E_3\bigr)
\]
が残る。許容データでは \(4/b^2\le1/2\) であり \(n\ge h/2>0\) である。
右辺の括弧の二倍は \(fa(\log b)^5\) 以下なので、\(n\) で割ればよい。
\leanlinks{\leanref{star_inequality}}
\end{proof}

\begin{corollary}[素数和の片側評価]\label{cor:prime-bound}
\(q_0\) を \(18480\) の正の倍数、\(d\equiv q_0/2-1\pmod{q_0}\) を素数とする。
任意の \(\varepsilon>0\) について、十分大きな実数 \(y\) に対し
\begin{equation}\label{eq:prime-bound}
 R_{q_0}(y)\le y^{-1/2+\varepsilon}
\end{equation}
が成り立つ。
\end{corollary}

\begin{proof}
\eqref{eq:main-ineq} を \eqref{eq:total-remainder} に代入し、
\(ma|\Phi_b|\) で割る。\(f=\varphi(q_0)|\Phi_a|\) および
\(|\Phi_a|/|\Phi_b|\le\log b\) を使うと、
\begin{equation}\label{eq:remainder-intermediate}
 R_{q_0}(x)
 \le\frac{\varphi(q_0)(\log b)^6}{m}
       +|R_{q_0}(x_1)|+\log b\,|R_1(x_1)|.
\end{equation}
\(0<\eta<1/100\)、\(7\eta\le\varepsilon\) とし、
命題 \ref{prop:admissible-data} を \(\eta\) で適用する。
十分大きな整数 \(x\) では
\(\log b\le2x^\eta\)、\(m\ge\sqrt{x}\)、\(\log x_1\ge x^{\eta/12}\) である。
補題 \ref{lem:fixed-modulus-pnt} に \(B\ge24/\eta\) を取れば、両方の
\(|R_{q_0}(x_1)|,|R_1(x_1)|\) は \(x^{-2}\) 以下となる。
したがって \eqref{eq:remainder-intermediate} の右辺は
\[
 64\varphi(q_0)x^{-1/2+6\eta}+x^{-2}+2x^{-2+\eta}
 \le (64\varphi(q_0)+3)x^{-1/2+6\eta}
\]
である。
\(x=\lfloor y\rfloor\) とすれば、補題 \ref{lem:fixed-modulus-pnt} の差の評価から
\(R_{q_0}(y)-R_{q_0}(x)\le26/y\) である。
\(\theta=-1/2+6\eta\) は \([-1,0]\) にあり、\(x\ge y/2\) なので
\(x^\theta\le2y^\theta\) である。
従って
\[
 R_{q_0}(y)\le (2(64\varphi(q_0)+3)+26)y^\theta .
\]
定数を十分大きな \(y^\eta\) に吸収し、\(7\eta\le\varepsilon\) を使えば
\eqref{eq:prime-bound} を得る。
\leanlinks{\leanref{eventually_primeRemainder_le_rpow}}
\end{proof}

\begin{proof}[定理 \ref{thm:grh} の証明]
\textup{(1)}
\(\chi\) を任意の法 \(q\) の指標とし、\(\RePart\rho>1/2\) とする。
\(\rho=1\) では主指標が除外され、非主指標については Dirichlet の
非消滅定理 \(L(1,\chi)\ne0\) による。以下では \(\rho\ne1\) とする。
\(q_0=18480q\)、\(d_0=q_0/2-1\) とおく。
\(d_0\equiv7\pmod8\) であり、\(q_0=2(d_0+1)\) かつ \(d_0\) は奇数なので
\((q_0,d_0)=1\) である。
\(1/2<\beta<\RePart\rho\) を選び、系 \ref{cor:prime-bound} を
\(\varepsilon=\beta-1/2\) で用いると、法 \(q_0\) の剰余類 \(d_0\) に属する
すべての素数 \(d\) について命題 \ref{prop:one-sided-criterion} の仮定が成立する。
\(\chi\) から誘導した法 \(q_0\) の指標を \(\chi_0\) とすれば
\(L(\rho,\chi_0)\ne0\) である。
誘導の公式
\[
 L(s,\chi_0)=L(s,\chi)
  \prod_{\substack{p\mid q_0\\p\nmid q}}\bigl(1-\chi(p)p^{-s}\bigr)
 \qquad(s\ne1)
\]
より \(L(\rho,\chi)\ne0\) を得る。

\textup{(2)}
\(0<\RePart s<1\) で \(L(s,\chi)=0\) とし、\(\chi\) を導く原始指標を \(\psi\) とする。
誘導に伴う有限個の Euler 因子はこの帯で非零なので \(L(s,\psi)=0\) である。
ガンマ因子はこの帯で正則かつ非零であるから、完備 \(L\) 関数も \(s\) で零となる。
\(\psi\) の関数等式により \(\psi^{-1}\) の完備 \(L\) 関数は \(1-s\) で零となり、
ガンマ因子を除いて \(L(1-s,\psi^{-1})=0\) を得る。
\(\psi^{-1}\) は \(\chi^{-1}\) を導くので \(L(1-s,\chi^{-1})=0\) となる。
\leanlinks{\leansrc{Lib/Analytic/ZeroSymmetry}{非原始指標を含む零点の反射}}
(1) を二つの零点にそれぞれ適用すると
\(\RePart s\le1/2\) および \(1-\RePart s\le1/2\) が従う。

\textup{(3)}
\(\zeta(s)=0\) とし、自明零点を除く。\(s=0\) は
\(\zeta(0)=-1/2\) により除かれ、\(s=1\) は極である。
\(\pi^{-s/2}\Gamma(s/2)\) はこの \(s\) で有限かつ非零なので、
完備ゼータ関数も \(s\) で零となる。
完備ゼータ関数は \(s\mapsto1-s\) で不変であり、
\(1-s\ne0,1\) でその値にガンマ因子の逆数を掛けると
\(\zeta(1-s)=0\) を得る。
法 \(1\) における (1) を二つの零点へ適用すると、
\(\RePart s=1/2\) が従う。
\leanlinks{\leanref{grh}}
\end{proof}
